\section{Comparação e entrega executiva}

\begin{frame}{Projetos do caderno — residenciais e comerciais}
\scriptsize
\begin{tabularx}{\textwidth}{c p{3.0cm} p{3.1cm} X}
\toprule
ID & Projeto & Diferencial de arquitetura & Foco elétrico \\
\midrule
R1 & Casa térrea & 2 suítes, banho social, quarto 3, garagem e hall interno & entrada BT, balanceamento e três chuveiros \\
R2 & Apartamento & 2 suítes, varanda, lavabo e serviço & unidade em EMUC, prumada e QD trifásico \\
C1 & Edifício corporativo & 4 pavimentos, 8 salas/tipo, WC M/F/PCD, elevador, escada & QGBT, elevador, SPDA, DPS \\
C2 & Centro empresarial & recepção, 10 salas, projetos, treinamento, copa e núcleo vertical & QGBT, TI, climatização, elevador e emergência \\
C3 & Loja de shopping & área de vendas, vitrines, provadores, caixa, estoque e sala elétrica & luminotécnica, teto refletido, cenas, alimentador e comissionamento \\
\bottomrule
\end{tabularx}
\end{frame}

\begin{frame}{Projetos do caderno — industriais}
\scriptsize
\begin{tabularx}{\textwidth}{c p{3.0cm} p{3.1cm} X}
\toprule
ID & Projeto & Diferencial de arquitetura & Foco elétrico \\
\midrule
I1 & Indústria leve & produção A/B, manutenção, almoxarifado, vestiários & motores/VFD/CCM com atendimento em BT \\
I2 & Indústria com SE & processo, embalagem, expedição, utilidades e refrigeração & MT, trafo 500 kVA, QGBT 1.000 A e SPDA \\
\bottomrule
\end{tabularx}
\begin{caixaProjeto}[title={Sete casos, quatro competências}]
Em todos os projetos, o aluno deve localizar e interpretar informações, dimensionar a solução, compreender a ligação/montagem e planejar execução, ensaios e entrega.
\end{caixaProjeto}
\end{frame}

\begin{frame}{O que muda de um projeto para outro}
\small
\begin{columns}[T,onlytextwidth]
\column{0.49\textwidth}
\begin{caixaProjeto}[title=Não são a mesma planta reciclada]
\begin{itemize}
\item geometrias e circulações diferentes;
\item usos e cargas próprios;
\item critérios luminotécnicos e cenas próprios na loja;
\item quadros em posições coerentes;
\item rotas compatíveis com cada arquitetura;
\item distribuição e unifilares diferentes.
\end{itemize}
\end{caixaProjeto}
\column{0.49\textwidth}
\begin{caixaNorma}[title=Norma conforme o caso]
\begin{itemize}
\item BT individual: NT.00001.EQTL;
\item múltiplas unidades: NT.00004.EQTL;
\item MT/subestação: NT.00002 e NT.00026;
\item SPDA: série ABNT NBR 5419;
\item loja: ABNT NBR ISO/CIE 8995-1, NBR 10898 e projeto de incêndio aprovado;
\item MT interna: ABNT NBR 14039.
\end{itemize}
\end{caixaNorma}
\end{columns}
\end{frame}

\begin{frame}{Checklist de projeto executivo}
\small
\begin{enumerate}
\item Usar a planta arquitetônica real, com cotas, níveis, layout e equipamentos.
\item Calcular iluminação/TUG e levantar TUE por dados de placa.
\item Dividir circuitos e definir fases.
\item Calcular $I_b$, escolher seção por $I_z$ e fatores de correção.
\item Verificar $I_b\le I_n\le I_z$ e condições de proteção.
\item Calcular queda de tensão, curto máximo/mínimo e seletividade.
\item Dimensionar N, PE, eletrodutos, bandejas, quadros e barramentos.
\item Integrar DR, DPS, aterramento e SPDA sem criar aterramentos independentes artificiais.
\item Em MT, solicitar dados de rede e coordenar proteção da Equatorial com a instalação interna.
\item Emitir memorial, documentação de responsabilidade técnica, comissionamento e as built.
\end{enumerate}
\end{frame}
